\section{Оптимальное управление}

В разделе 3 была получена допустимая область параметров регулятора, в которой одновременно выполняются условия устойчивости системы и ограничения
\begin{equation}
|M_{\text{м}}|_{\max} < 150 \text{ Нм},
\qquad
|u|_{\max} < 220 \text{ В}.
\end{equation}

На расчётной сетке были получены значения
\begin{equation}
|M_{\text{м}}|_{\max} \approx 149.99 \text{ Нм},
\qquad
|u|_{\max} \approx 62.60 \text{ В},
\end{equation}
а также
\begin{equation}
t_{\text{п,min}} \approx 1.1 \text{ с},
\qquad
t_{\text{п,max}} \approx 40 \text{ с}.
\end{equation}

Значение \(t_{\text{п,max}} = 40\) с соответствует точкам, в которых переходный процесс не вошёл в десятипроцентную полосу на интервале моделирования.

Теперь внутри допустимой области необходимо выбрать параметры регулятора, обеспечивающие наилучшее качество переходного процесса. Для этого вводится интегральный показатель качества
\begin{equation}
J_{\text{оп}}(\rho)
=
\int_0^{t_{\text{п}}}
\left(
\psi_{\text{м}}^2
+
\rho u^2
\right)dt,
\end{equation}
где
\begin{equation}
\psi_{\text{м}}(t)
=
\varphi_{\text{м}}(t)-\varphi^*
=
\varphi_{\text{м}}(t)-\pi.
\end{equation}

Параметр \(\rho\) задаёт вес управляющего сигнала в функционале качества.

При \(\rho=0\) учитывается только ошибка регулирования. При увеличении \(\rho\) возрастает влияние управляющего сигнала.

Расчёт проводился на той же сетке параметров, что и в разделе 3. В расчёт принимались только точки допустимой области.

Для каждой точки сначала определялось время переходного процесса \(t_{\text{п}}\), после чего вычислялись интегралы
\begin{equation}
J_{\psi}
=
\int_0^{t_{\text{п}}}
\psi_{\text{м}}^2(t)\,dt,
\qquad
J_u
=
\int_0^{t_{\text{п}}}
u^2(t)\,dt.
\end{equation}

Численное интегрирование выполнялось по временным отсчётам. При шаге интегрирования \(\Delta t\)
\begin{equation}
J_{\psi}
\approx
\sum_{j=0}^{N}
\psi_{\text{м}}^2(t_j)\Delta t,
\qquad
J_u
\approx
\sum_{j=0}^{N}
u^2(t_j)\Delta t,
\end{equation}
где
\begin{equation}
t_N=t_{\text{п}}.
\end{equation}

После этого вычислялся функционал
\begin{equation}
J_{\text{оп}}
=
J_{\psi}
+
\rho J_u.
\end{equation}
